\section{Implementasi Program}
\subsection{Lingkungan dan Struktur Program}
Algoritma local search diimplementasikan dengan Python sesuai spesifikasi.
Teknologi visualisasi dibebaskan. Repository wajib memiliki folder
\texttt{src/}, \texttt{docs/}, \texttt{output/}, dan \texttt{README.md}.
Laporan PDF ditempatkan dalam \texttt{docs/}, sedangkan hasil run disimpan
dalam \texttt{output/}~\cite{spesifikasi}.


\subsection{Representasi State dan Pemeriksaan Validitas}
Komponen representasi menghubungkan atribut paket, posisi, orientasi, status
di dalam/luar, dan keadaan ruang truk. Validator harus menerapkan aturan pada
Bagian~\ref{sec:validasi} secara konsisten pada inisialisasi maupun perubahan state.


\subsection{Initial State, Neighbor, dan Fungsi Objektif}
Komponen ini menyediakan konfigurasi acak, kandidat dari tiga jenis move,
dan perhitungan nilai muatan. Apabila digunakan cache atau evaluasi
inkremental, hasilnya perlu diperiksa terhadap perhitungan ulang seluruh state.


\subsection{Implementasi Hill-Climbing}


\subsection{Implementasi Simulated Annealing}


\subsection{Implementasi Genetic Algorithm}


\subsection{Visualisasi dan Penyimpanan Hasil}
Visualisasi harus menampilkan keadaan seluruh paket dan peti kemas, bukan
hanya nilai agregat. Paket di luar dapat ditunjukkan pada panel terpisah
selama identitas dan statusnya tetap jelas. Output eksperimen perlu mengaitkan
state, histori, durasi, konfigurasi, dan seed dengan run yang sama.
